%! Factoring Special Forms — Unit 1, Lesson 1.5 — Homework
\documentclass[10pt]{article}
\usepackage{atda-article}
\usepackage{atda-boxes}
\newcommand{\TallMath}[1]{$\displaystyle #1\rule[-1.4em]{0pt}{3.2em}$}
\setlength{\workrowsep}{8pt}

\begin{document}

\pageheader{Unit 1, Lesson 1.5}{Homework}
% NAMESTRIP: no \namedateperiod here. The student writes their name once, on the
% cover this component is stapled behind. See references/conventions.md.

\vspace{0.15in}

\boxguard[16]
\begin{notesbox}{Practice --- The Three Special Forms}
\small
Show your work. \textbf{Name the pattern} before you factor, factor \textbf{completely} every time,
and check by multiplying back. If nothing factors over the integers, write \emph{prime}.

\begin{enumerate}[leftmargin=1.6em, itemsep=6pt, topsep=4pt]

\item Factor: $x^2-81$
\begin{work}
  x^2-81 &= x^2-9^2 \\
         &= (x+9)(x-9)
\end{work}

\item Factor: $x^2+16x+64$
\begin{work}
  2(x)(8) &= 16x \quad\checkmark \\
  x^2+16x+64 &= (x+8)^2
\end{work}

\item Factor: $x^2-20x+100$
\begin{work}
  2(x)(10) &= 20x \quad\checkmark \\
  x^2-20x+100 &= (x-10)^2
\end{work}

\item Factor: $49x^2-16$
\begin{work}
  49x^2-16 &= (7x)^2-4^2 \\
           &= (7x+4)(7x-4)
\end{work}

\item Factor: $x^2+36$
\begin{work}
  x^2+36 &= x^2+6^2, \quad \text{a \emph{sum} of squares} \\
  x^2+36 &\text{ is prime --- no integer pair has product } 36 \text{ and sum } 0.
\end{work}

\item Factor: $x^3+125$
\begin{work}
  x^3+125 &= x^3+5^3 \\
          &= (x+5)\left(x^2-5x+25\right)
\end{work}

\item Factor: $x^3-27$
\begin{work}
  x^3-27 &= x^3-3^3 \\
         &= (x-3)\left(x^2+3x+9\right)
\end{work}

\item Factor: $8x^3+1$
\begin{work}
  8x^3+1 &= (2x)^3+1^3 \\
         &= (2x+1)\left(4x^2-2x+1\right)
\end{work}

\item Factor completely --- GCF first, then the pattern: $3x^3-75x$
\begin{work}
  3x^3-75x &= 3x\left(x^2-25\right) \\
           &= 3x(x+5)(x-5)
\end{work}

\end{enumerate}
\end{notesbox}

\vspace{0.10in}

\boxguard[16]
\begin{scenariobox}[In Context --- The Courtyard Garden]{royal}
\small
The horticulture class is planting a \textbf{square} courtyard garden with a \textbf{square} stone
fountain at its center. The plan sheet gives only the planted area:
\[ A(x) = 25x^2-4 \quad \text{square feet,} \]
which is the garden's $5x$ by $5x$ footprint with the fountain's $2$ ft by $2$ ft base removed.

\begin{enumerate}[leftmargin=1.6em, itemsep=6pt, topsep=4pt, start=10]

\item Factor $A(x)$ to get the dimensions of the rectangle the planted region rearranges into.
\begin{work}
  25x^2-4 &= (5x)^2-2^2 \\
          &= (5x+2)(5x-2)
\end{work}

\item The class settles on $x=3$. Give both dimensions and the planted area, in feet, with units ---
then check the area against the original polynomial.
\begin{work}
  5x+2 &= 17 \text{ ft}, \qquad 5x-2 = 13 \text{ ft} \\
  17 \cdot 13 &= 221 \text{ sq ft} \\
  25(9)-4 &= 221 \text{ sq ft} \quad\checkmark
\end{work}

\item The class also has to buy border fencing. In one sentence, say what the factored form tells
them that the expanded form does not.

\writeline

\end{enumerate}
\end{scenariobox}

\vspace{0.10in}

\boxguard[16]
\begin{scenariobox}[A First Look Ahead]{goldacc}
\small
\begin{enumerate}[leftmargin=1.6em, itemsep=6pt, topsep=4pt, start=13]

\item A courtyard patio is a square $x$ feet on a side with a square planter $6$ feet on a side at
its center, so the paved area is $A(x) = x^2-36$ square feet.
\begin{enumerate}[label=(\alph*), leftmargin=1.6em, itemsep=4pt, topsep=3pt]
  \item Factor $A(x)$.
\begin{work}
  x^2-36 &= x^2-6^2 \\
         &= (x+6)(x-6)
\end{work}
  \item For what value of $x$ is the paved area \textbf{zero} --- that is, when does the planter
        fill the whole patio? Use the factored form, and the fact that a product of two numbers is
        zero only when one of the two \emph{is} zero.
\begin{work}
  x+6 &= 0 \quad\Rightarrow\quad x = -6 \; \text{(rejected --- a length cannot be negative)} \\
  x-6 &= 0 \quad\Rightarrow\quad x = 6 \\
  6^2-36 &= 0 \quad\checkmark
\end{work}
\end{enumerate}

\end{enumerate}
\end{scenariobox}

\vspace{0.10in}

\boxguard[18]
\begin{extensionbox}
\small
\textbf{Two variables, same pattern.} A landscaper writes a plot's area as $4x^2-9y^2$ square feet,
where $x$ and $y$ are both measured in feet. Factor it, then multiply your answer back out to
\textbf{verify the identity}. What are $a$ and $b$ here?
\begin{work}
  4x^2-9y^2 &= (2x)^2-(3y)^2 \\
            &= (2x+3y)(2x-3y) \\
  (2x+3y)(2x-3y) &= 4x^2-6xy+6xy-9y^2 \\
                 &= 4x^2-9y^2 \quad\checkmark
\end{work}
\writeline
\end{extensionbox}

\vspace{0.10in}

\begin{remindbox}
\small
\textbf{Next: Lesson 1.6 --- Solving Polynomial Equations by Factoring}
(\texttt{A2.EI.2b}, \texttt{A2.EI.6a--b}). Item 13(b) already ran the whole method once: you
factored, then asked which value of $x$ makes a factor zero. In 1.6 that move gets its name --- the
\textbf{zero product property} --- and it turns every factoring problem in this unit into a solved
equation.
\end{remindbox}

\end{document}
